\section{Methods}
\label{sec:methods}

\subsection{Approach}
\label{sec:approach}

Model-free reinforcement learning treats the simulator as a black box and
estimates the policy gradient from sampled returns alone. This is robust and
makes almost no assumptions, but the estimate is high-variance and
sample-hungry, and on flow-control problems it can take tens of thousands of
environment interactions to converge. When the simulator is itself
differentiable, a second route becomes available: differentiate the
accumulated reward \emph{through} the solver and obtain the policy gradient
from the chain rule directly. This \emph{analytic} (or first-order) gradient
is far lower in variance and, on smooth control problems, converges in a
fraction of the interactions. It is the mechanism behind
differentiable-physics methods such as SHAC~\cite{shac2022},
AHAC~\cite{ahac2024} and SAPO~\cite{sapo2025} in robotics, behind
differentiable-CFD control with JAX-CFD~\cite{jaxcfd2021} and
JAX-Fluids~\cite{jaxfluids2022}, and behind the GPPO result on Kolmogorov
flow that this paper follows up~\cite{hydrogym2026}.

The derivative that makes this route fast also makes it fragile.
Backpropagating through a long or stiff rollout produces gradients that
explode or vanish, and the estimate is biased whenever the horizon is
truncated or the dynamics are contact-rich or
chaotic~\cite{metz2021,suh2022}. The established mitigations are to use a
\emph{short} differentiation horizon, to forward-propagate only the
deterministic mean action so the differentiated chain stays smooth, and to
\emph{mix} the analytic gradient with the ordinary model-free gradient rather
than trust it alone~\cite{shac2022,suh2022,gippo2023}. Chaotic flow is
exactly where this fragility is expected to matter most: neighbouring
trajectories separate exponentially, so the through-solver gradient of a
many-step rollout can be swamped by noise, and the predictability horizon
shrinks as the Reynolds number rises~\cite{metz2021,boffetta2017}. Whether
the analytic gradient still helps in that regime --- and whether an
off-policy replay buffer absorbs the instability better than an on-policy
update --- is the question this study is built to answer.

The rest of this section constructs the four algorithms used to test it. We
fix a single environment and control problem (\cref{sec:env}) and cast it as
a Markov decision process (\cref{sec:mdp}); we state the two model-free
estimators, PPO and SAC (\cref{sec:ppo,sec:sac}); we then define the analytic
gradient and the fixed 50/50 mixing rule (\cref{sec:analytic}); and finally
we assemble the on-policy variant GPPO (\cref{sec:gppo}) and the off-policy
variant GSAC (\cref{sec:gsac}). The factorial design that crosses these
choices is described first.

\subsection{Overall design}
\label{sec:design}

We compare four algorithms in a $2\times 2$ factorial design that crosses the
RL family (on-policy vs.\ off-policy) with analytic-gradient injection
(absent vs.\ present), so that each factor and their interaction can be read
off directly (\cref{fig:design}). All four share one environment, one network
architecture, one optimizer, and one training budget (\cref{sec:setup}); they
are evaluated at a regular ($\Re=45$) and a chaotic ($\Re=300$) operating
point of the same Kolmogorov-flow solver.

\begin{figure*}[t]
\centering
\begin{tikzpicture}[node distance=6mm]
  \node[font=\small\bfseries] (c1) at (0,2.1) {model-free gradient};
  \node[font=\small\bfseries] (c2) at (4.2,2.1) {$+$ analytic gradient};
  \node[font=\small\bfseries,rotate=90] at (-2.6,1.0) {on-policy};
  \node[font=\small\bfseries,rotate=90] at (-2.6,-1.4) {off-policy};

  \node[block,minimum width=30mm] (ppo)  at (0,1.0)  {\textbf{PPO}\\clipped surrogate\\ \cite{ppo2017}};
  \node[anblock,minimum width=30mm] (gppo) at (4.2,1.0) {\textbf{GPPO}\\PPO $+$ through-sim.\\reward gradient};
  \node[block,minimum width=30mm] (sac)  at (0,-1.4) {\textbf{SAC}\\max-entropy, twin\\critics \cite{sac2018}};
  \node[anblock,minimum width=30mm] (gsac) at (4.2,-1.4) {\textbf{GSAC}\\SAC $+$ through-sim.\\reward gradient (new)};

  \draw[flow] (ppo) -- (gppo);
  \draw[flow] (sac) -- (gsac);
  \draw[flow] (ppo) -- (sac);
  \draw[flow] (gppo) -- (gsac);

  % shared environment, centred under both columns; straight vertical edges
  \node[envblock,minimum width=76mm,anchor=north] (env) at (2.1,-3.5)
    {\textbf{Shared environment} --- real \texttt{hydrogym.jax} Kolmogorov flow\\
    $64\times 64$ grid (4096 DOF), pseudo-spectral IMEX-RK4CN solver\\
    regime knob: $\Re \in \{45\ (\text{regular}),\ 300\ (\text{chaotic})\}$};
  \draw[flow] (sac.south) -- (sac.south |- env.north);
  \draw[flow] (gsac.south) -- (gsac.south |- env.north);
\end{tikzpicture}
\caption{The $2\times 2$ experimental design. Horizontal edges add the
analytic (backpropagated-through-the-simulator) reward gradient; vertical
edges switch the RL family. GSAC (bottom right) is introduced here. All four
algorithms are run at both Reynolds numbers with identical budgets.}
\label{fig:design}
\end{figure*}

\subsection{Environment and control problem}
\label{sec:env}

Experiments run on the real
\texttt{hydrogym.jax.\allowbreak envs.\allowbreak kolmogorov.\allowbreak KolmogorovFlow} environment: a
$64\times 64$ grid (4096 degrees of freedom), pseudo-spectral IMEX-RK4CN
time integration~\cite{hydrogym2026,jaxcfd2021}, with four low-wavenumber
sinusoidal forcing modes as actuators (\cref{fig:env}). The 2D
incompressible Navier--Stokes equations in vorticity form are integrated as
$\omega_{t+1}=f_{\mathrm{sim}}(\omega_t, a_t;\Re)$, where the discrete
operator $f_{\mathrm{sim}}$ is written in JAX~\cite{jax2018} and is differentiable
end-to-end.

\begin{figure*}[t]
\centering
\begin{tikzpicture}[node distance=6mm]
  % flow domain
  \draw[thick] (0,0) rectangle (2.4,2.4);
  \foreach \y in {0.4,0.8,1.2,1.6,2.0}{\draw[gray!55] (0,\y) -- (2.4,\y);}
  \node[font=\scriptsize] at (1.2,2.75) {$64\times 64$ vorticity field $\omega_t$};
  % actuator arrows (sinusoidal forcing modes)
  \foreach \y/\s in {0.4/0.3,0.8/-0.3,1.2/0.3,1.6/-0.3,2.0/0.3}{
    \draw[-{Latex[length=1.4mm]},blue!70!black] (0.12,\y) -- ($(0.12,\y)+(\s,0)$);
    \draw[-{Latex[length=1.4mm]},blue!70!black] (2.28,\y) -- ($(2.28,\y)-(\s,0)$);}
  \node[blue!70!black,font=\scriptsize] at (1.2,-0.55)
    {4 forcing modes $a_t\in\mathbb{R}^4$};

  % observation
  \node[databox,align=center] (obs) at (6.1,1.9)
    {\scriptsize $8\times 8$ subsampled\\ velocity magnitude\\ $s_t\in\mathbb{R}^{64}$};
  \draw[flow] (2.4,1.7) -- (obs.west) node[midway,above,font=\scriptsize]{observe};

  % reward
  \node[databox,align=center] (rew) at (6.1,0.35)
    {\scriptsize $r_t = -\big(E(\omega_t)+c\,\lVert a_t\rVert^2\big)$\\[1pt]
     (flow energy $+$ control cost)};
  \draw[flow] (obs.south) -- (rew.north);
\end{tikzpicture}
\caption{Kolmogorov-flow control problem. The policy observes an $8\times 8$
subsampled velocity magnitude ($s_t\in\mathbb{R}^{64}$), commands four
sinusoidal forcing-mode amplitudes ($a_t\in\mathbb{R}^{4}$), and is rewarded
for reducing flow energy at bounded control cost. The library default
$\alpha_{\text{reward}}=0$ (control-cost only) is overridden to
$\alpha_{\text{reward}}=1$ so that flow energy enters the reward.}
\label{fig:env}
\end{figure*}

Two library quirks required a subclass. (i) The environment step ignores the
per-call time-step and action-duration arguments and re-reads its own
\texttt{default\_params}, so we override that property on a subclass
(\texttt{FastKolmogorovFlow}) to reach proof-of-concept scale.
(ii) The library default $\alpha_{\text{reward}}=0$ makes the reward
control-cost only; we set $\alpha_{\text{reward}}=1$ (\cref{fig:env}).

At this grid/$\dd t$ combination ($64\times 64$, $\dd t=0.001$) the solver was
empirically observed to diverge numerically near $\Re\approx 350$--$380$ ---
even under zero action from purely analytic initial conditions, and not
resolved by reducing $\dd t$ tenfold (so not a simple CFL issue). We
therefore use $\Re=300$ for the chaotic regime --- the highest value that
stays below this limit while forming a meaningful contrast with $\Re=45$ ---
rather than the originally intended $\Re=450$.

\subsection{Markov decision process}
\label{sec:mdp}

The control task is a Markov decision process
$(\mathcal{S},\mathcal{A},P,r,\gamma)$ with
$s_t\in\mathbb{R}^{64}$, $a_t\in\mathbb{R}^{4}$, deterministic transition
$s_{t+1}=g\big(f_{\mathrm{sim}}(\omega_t,a_t;\Re)\big)$ ($g$ the observation
map), reward as in \cref{fig:env}, and discount $\gamma=0.99$. An episode is
$T=15$ environment steps from a fixed reset state $s_0$. The policy
$\pi_\theta(a\mid s)$ and the value / $Q$ networks are 2-hidden-layer MLPs
(64 units, $\tanh$). The objective is the discounted return
\begin{equation}
  J(\theta)=\mathbb{E}_{\tau\sim\pi_\theta}\!\left[\textstyle\sum_{t=0}^{T-1}\gamma^{t}\,r(s_t,a_t)\right].
  \label{eq:return}
\end{equation}
All four algorithms optimize \cref{eq:return}; they differ only in the
gradient estimator they apply to it.

\subsection{PPO: on-policy baseline}
\label{sec:ppo}

PPO is an on-policy policy-gradient method: it collects one rollout under
the current policy, estimates each action's advantage along it, and takes a
single trust-region-style step that is clipped to stay close to the policy
that generated the data. The rollout is then discarded, so PPO consumes
each unit of environment experience through exactly one gradient update.

The clipped surrogate is due to Schulman et al.\ as a first-order
approximation to the TRPO trust region~\cite{ppo2017,trpo2015}, and the
Generalized Advantage Estimator (GAE) used to form the advantage is from the
same group~\cite{gae2015}. All gradient steps in this study, for every
algorithm, use Adam~\cite{kingma2015adam}, as in the original PPO and
HydroGym~\cite{hydrogym2026} implementations.

The policy is a diagonal Gaussian: the MLP produces a bounded mean
$a_{\mu}(s)=0.5\tanh(\cdot)$ and a state-independent learnable
$\log\sigma$; the executed action is $a\sim\mathcal{N}(a_\mu(s),\sigma^2)$
and is \emph{not} re-clipped, so sampled $\lvert a\rvert$ can exceed
$0.5$ (observed up to $\approx 1.2$ early in training). Advantages use
GAE~\cite{gae2015}:
\begin{equation}
  \hat{A}_t=\sum_{l=0}^{T-t-1}(\gamma\lambda)^{l}\delta_{t+l},\qquad
  \delta_t=r_t+\gamma V(s_{t+1})-V(s_t),
  \label{eq:gae}
\end{equation}
with $\lambda=0.95$. PPO maximizes the clipped
surrogate~\cite{ppo2017}, derived from the TRPO trust
region~\cite{trpo2015}:
\begin{equation}
  \begin{split}
    \mathcal{L}^{\mathrm{PPO}}(\theta)&=\mathbb{E}_t\!\Big[\min\!\big(\rho_t(\theta)\hat{A}_t,\\
    &\qquad \mathrm{clip}(\rho_t(\theta),1-\epsilon,1+\epsilon)\hat{A}_t\big)\Big],\\
    \rho_t(\theta)&=\frac{\pi_\theta(a_t\mid s_t)}{\pi_{\theta_{\text{old}}}(a_t\mid s_t)},
  \end{split}
  \label{eq:ppo}
\end{equation}
with $\epsilon=0.2$. The total loss adds a value term and an entropy bonus,
\begin{equation}
  \begin{split}
    \mathcal{L}&=-\mathcal{L}^{\mathrm{PPO}}+0.5\,\mathcal{L}^{V}-0.001\,\mathcal{H}[\pi_\theta],\\
    \mathcal{L}^{V}&=\mathbb{E}_t\big[(V_\phi(s_t)-\hat{R}_t)^2\big].
  \end{split}
  \label{eq:ppoloss}
\end{equation}

Our loop performs \emph{one
full-batch gradient step per episode} on \cref{eq:ppoloss} --- no PPO
epochs, no minibatching. This is a deliberately minimal PPO: it keeps the
on- and off-policy loops comparable and the analytic-gradient mix simple,
but ``PPO'' here is weaker than a fully tuned implementation. The comparison
of interest is relative (PPO vs.\ GPPO, SAC vs.\ GSAC), and all four share
this budget.

We use PPO as the on-policy baseline, rather than a
different on-policy algorithm, because it is exactly the update the
HydroGym paper's GPPO result is built on top of~\cite{hydrogym2026}:
answering whether that result survives a chaotic regime requires holding
the base update fixed and changing only the analytic-gradient mixing
(\cref{sec:gppo}), not also switching algorithms. Its clipped surrogate also
gives a clean model-free-only control arm for the on-policy row of the
$2\times 2$ design (\cref{fig:design}).

\subsection{SAC: off-policy baseline}
\label{sec:sac}

SAC is an off-policy, maximum-entropy actor-critic method: transitions are
stored in a replay buffer and resampled many times, so one environment step
contributes to many gradient updates instead of one. Its actor is a
squashed Gaussian whose action is drawn by reparameterized sampling, so the
entropy-regularized objective can be differentiated through the sampled
action directly rather than estimated by a likelihood-ratio term.

SAC and its maximum-entropy objective are due to Haarnoja et
al.~\cite{sac2018,sacapps2018}, building on Soft
Q-learning~\cite{softq2017}; its twin critics adopt the overestimation-bias
fix introduced for off-policy actor-critic methods by
TD3~\cite{fujimoto2018td3}, and the reparameterized actor gradient traces to
the stochastic backpropagation of Rezende et
al.~\cite{rezende2014stochastic}, the same mechanism used earlier for
deterministic actors in DDPG~\cite{lillicrap2016ddpg}. SAC uses twin
$Q$-critics $Q_{\psi_1},Q_{\psi_2}$ and target networks to optimize:
\begin{equation}
  J_{\mathrm{ME}}(\pi)=\mathbb{E}\!\left[\textstyle\sum_t r(s_t,a_t)+\alpha\,\mathcal{H}\!\big(\pi(\cdot\mid s_t)\big)\right],
  \label{eq:sacobj}
\end{equation}
with fixed temperature $\alpha=0.1$ (not auto-tuned). The actor and critic
losses are
\begin{align}
  J_\pi(\varphi)     &=\mathbb{E}_{s\sim\mathcal{D},\,a\sim\pi_\varphi}\!\big[\alpha\log\pi_\varphi(a\mid s)-\min_i Q_{\psi_i}(s,a)\big],
  \label{eq:sacactor}\\
  J_Q(\psi_i)        &=\mathbb{E}_{(s,a,r,s')\sim\mathcal{D}}\!\Big[\big(Q_{\psi_i}(s,a)-y\big)^2\Big],\\
  y&=r+\gamma\big(\min_j Q_{\bar\psi_j}(s',a')-\alpha\log\pi_\varphi(a'\mid s')\big),
  \label{eq:saccritic}
\end{align}
with soft target updates
$\bar\psi_i\leftarrow(1-\tau_s)\bar\psi_i+\tau_s\psi_i$, $\tau_s=0.005$. Per
episode, the rollout's 15 transitions are added to a replay buffer $\mathcal{D}$
(capacity $=n_{\text{ep}}\times 15$) and 10 gradient updates are performed
with batch size 64.

SAC is the natural off-policy counterpart to PPO
along exactly the factor this study varies (\cref{fig:design}): both
optimize a stochastic policy with an actor-critic architecture, so replacing
PPO's rollout-and-discard loop with SAC's replay buffer isolates the
on-/off-policy factor without also changing the policy family, as switching
to a deterministic actor (DDPG/TD3) would. Its replay buffer also gives the
off-policy family a mechanism, buffer dilution, that PPO structurally lacks;
\cref{sec:buffer} uses this to explain why on-policy and off-policy runs can
behave differently immediately after a chaotic-regime trigger episode even
when neither family is more failure-prone on average. We sample the buffer
uniformly rather than by TD-error priority~\cite{schaul2016per}, so that any
such difference is attributable to the on-/off-policy structure itself and
not to a secondary sampling scheme.

\subsection{Analytic policy gradient}
\label{sec:analytic}

The analytic gradient is the derivative of a short rollout's accumulated
reward with respect to the policy parameters, obtained by unrolling the
policy's \emph{deterministic mean} action through the differentiable solver
and differentiating end-to-end with reverse-mode automatic differentiation.
It is a low-variance but biased estimator, so its length is capped at
$H=6$ steps and it is mixed with the model-free gradient at a fixed 50/50
weight rather than trusted alone (\cref{fig:analytic}).

Backpropagating through a rollout is the classical
backpropagation-through-time construction~\cite{werbos1990bptt}; SHAC~\cite{shac2022}
and AHAC~\cite{ahac2024} apply the same short-horizon truncation to
differentiable simulators, and Metz et al.~\cite{metz2021} and Parmas et
al.~\cite{parmas2018pipps} show that such backpropagated gradients can
explode and lose direction under chaotic dynamics. Over $H=6$ steps we
forward-propagate the policy's mean action (no sampling, so the whole chain
stays smooth) and differentiate the negative accumulated reward:
\begin{equation}
  \begin{split}
    \Lan(\theta)&=-\sum_{t=1}^{H}r\big(s_t,a_\mu(s_t;\theta)\big),\\
    s_t&=f_{\mathrm{sim}}(s_{t-1},a_\mu(s_{t-1};\theta);\Re),\\
    \gan&=\nabla_\theta\Lan(\theta),
  \end{split}
  \label{eq:analytic}
\end{equation}
computed by \texttt{jax.grad}. The final gradient mixes $\gan$ with the
model-free gradient $g_{\mathrm{mf}}\in\{g_{\mathrm{PPO}},g_\pi\}$ at equal
weight:
\begin{equation}
  g_{\mathrm{final}}=\beta\,\gan+(1-\beta)\,g_{\mathrm{mf}},\qquad \beta=0.5,
  \label{eq:mix}
\end{equation}
the $\alpha=0.5$ special case of the $\alpha$-graded estimator of Suh et
al.~\cite{suh2022}, of the total-gradient decomposition of Parmas and
Sugiyama~\cite{parmas2018total}, and the non-adaptive limit of the GI-PPO
$\alpha$-policy~\cite{gippo2023}. \Cref{eq:mix} is the only structural
difference between PPO and GPPO, and between SAC and GSAC.

\begin{figure*}[t]
\centering
\resizebox{\textwidth}{!}{%
\begin{tikzpicture}
  % forward rollout, row 1
  \node[block] (pi0) at (0,0) {$\pi_\theta$};
  \node[block] (a1)  at (2.5,0) {$a_\mu(s_0)$};
  \node[envblock] (f1) at (5.0,0) {$f_{\mathrm{sim}}$};
  \node[block] (s1)  at (7.5,0) {$s_1$};
  \node[databox] (r1) at (10.0,0) {$r_1$};
  % forward rollout, row 2
  \node[block] (a2)  at (7.5,-2.5) {$a_\mu(s_1)$};
  \node[envblock] (f2) at (10.0,-2.5) {$f_{\mathrm{sim}}$};
  \node[font=\small] (dots) at (12.0,-2.5) {$\cdots$};
  \node[databox] (rH) at (13.6,-2.5) {$r_H$};
  \node[anblock] (L) at (13.6,-5.0) {$\displaystyle \Lan=-\sum_{t=1}^{H}r_t$};

  \draw[flow] (pi0)--(a1); \draw[flow] (a1)--(f1); \draw[flow] (f1)--(s1);
  \draw[flow] (s1)--(r1);
  \draw[flow] (s1)--(a2);
  \draw[flow] (a2)--(f2); \draw[flow] (f2)--(dots); \draw[flow] (dots)--(rH);
  \draw[flow] (rH)--(L);
  \draw[flow] (r1.east) -- (15.0,0) -- (15.0,-5.0) -- (L.east);

  % backward analytic gradient: one bus under the rollout, orthogonal drops
  \draw[grad] (L.west) -- (2.5,-5.0) -- (a1.south);
  \draw[grad] (5.0,-5.0) -- (f1.south);
  \draw[grad] (7.5,-5.0) -- (a2.south);
  \draw[grad] (10.0,-5.0) -- (f2.south);
  \node[font=\scriptsize,text=red!70!black,anchor=north] at (6.2,-5.15)
    {$\nabla_\theta$ through solver \& policy};

  % mixer and optimiser
  \node[block,fill=yellow!12] (mix) at (13.6,-7.4)
    {$g_{\mathrm{final}}=\tfrac12\gan+\tfrac12 g_{\mathrm{mf}}$};
  \node[block] (mf) at (8.4,-7.4) {$g_{\mathrm{mf}}$\\(PPO / SAC actor)};
  \node[block] (adam) at (18.0,-7.4) {Adam\\ $\theta\!\leftarrow\!\theta-\eta g_{\mathrm{final}}$};
  \draw[grad] (L) -- (mix);
  \draw[flow] (mf) -- (mix);
  \draw[flow] (mix) -- (adam);
  \draw[flow] (adam.south) -- (18.0,-9.2) -- (0,-9.2) -- (pi0.south);
\end{tikzpicture}}
\caption{Analytic-gradient computation graph (\cref{eq:analytic,eq:mix}).
Solid arrows: forward rollout of the \emph{deterministic mean} action
through the differentiable solver for $H=6$ steps. Dashed red arrows:
reverse-mode differentiation of the accumulated reward $\Lan$ with respect to
$\theta$ via \texttt{jax.grad}. The analytic gradient is averaged 50/50 with
the model-free gradient and applied with Adam. Replacing the mixer with
$g_{\mathrm{final}}=g_{\mathrm{mf}}$ recovers plain PPO / SAC.}
\label{fig:analytic}
\end{figure*}

We fix $\beta=0.5$ and $H=6$
rather than adapting either online, because this study asks whether
analytic-gradient injection helps or hurts as the regime turns chaotic, and
an adaptive scheme such as GI-PPO's $\alpha$-policy or SAPO's entropy-based
smoothing would let the algorithm partially opt out of the analytic term
exactly when the chaotic regime makes it unreliable --- confounding the
effect under test. \Cref{sec:gppo,sec:gsac} apply this same fixed rule to
the on-policy and off-policy actor gradient respectively, so that GPPO and
GSAC differ from PPO and SAC by exactly one mechanism.

\subsection{GPPO: on-policy variant}
\label{sec:gppo}

GPPO adds the analytic-gradient mix of \cref{eq:mix} to the PPO update: from
the same fixed reset state, a deterministic $H=6$-step analytic rollout runs
alongside the standard stochastic on-policy rollout, and the two resulting
gradients are averaged before a single Adam step (\cref{fig:gppo,alg:gppo}).
Every other setting --- rollout length, GAE, the clip range, network
architecture, and learning rate --- is identical to plain PPO (\cref{sec:ppo}).

This is the same mechanism as the HydroGym paper's GPPO, reported to cut
training iterations by at least \SI{65}{\percent} over PPO on Kolmogorov
flow~\cite{hydrogym2026}. That paper does not report its horizon length or
mixing weight, so ours follow the general short-horizon-and-mix recipe of
SHAC, AHAC, and GI-PPO (\cref{sec:analytic}) rather than an attempt to match
unpublished hyperparameters.

\begin{figure*}[t]
\centering
\resizebox{\textwidth}{!}{%
\begin{tikzpicture}[node distance=6mm and 10mm]
  \node[block] (reset) {reset $\to s_0$};
  \node[block,right=of reset] (roll) {stochastic rollout\\ $T=15$, $a\!\sim\!\mathcal{N}$};
  \node[block,right=of roll] (gae) {GAE \eqref{eq:gae}\\ $\hat A_t,\hat R_t$};
  \node[block,right=of gae] (gppo) {$g_{\mathrm{PPO}}=\nabla_\theta\mathcal{L}$\\ \eqref{eq:ppoloss}, clip $\epsilon{=}0.2$};
  \node[block,right=of gppo,fill=yellow!12] (mix) {mix \eqref{eq:mix}\\ $\beta=0.5$};
  \node[block,right=of mix] (adam) {single Adam step\\ $\eta=3\!\times\!10^{-4}$};
  \node[anblock,below=10mm of gae] (an) {analytic branch\\ $H{=}6$, mean action\\ $\gan$ \eqref{eq:analytic}};
  \draw[flow] (reset)--(roll); \draw[flow] (roll)--(gae); \draw[flow] (gae)--(gppo);
  \draw[flow] (gppo)--(mix); \draw[flow] (mix)--(adam);
  \draw[flow] (reset.south) |- (an.west);
  \draw[grad] (an.east) -| (mix.south);
  % next episode: below everything, re-enter reset from the left
  \coordinate (low) at ($(an.south)+(0,-8mm)$);
  \draw[flow] (adam.south) -- (adam.south |- low) -- ($(reset.west |- low)+(-6mm,0)$)
      -- ($(reset.west)+(-6mm,0)$) -- (reset.west);
  \node[font=\scriptsize,above,anchor=south] at ($(reset.west |- low)+(2cm,0)$) {next episode};
\end{tikzpicture}}
\caption{GPPO training loop (one episode). The analytic branch is a
deterministic 6-step unroll that runs in parallel with the standard on-policy
update; the two gradients are averaged before a single Adam step. Removing
the red branch gives PPO.}
\label{fig:gppo}
\end{figure*}

\begin{algorithm}[t]
\caption{GPPO --- one training episode ($\beta=0.5$, $H=6$)}
\label{alg:gppo}
\DontPrintSemicolon
\KwIn{policy $\theta$, value $\phi$, optimizer state, fixed reset state $s_0$}
$\tau \leftarrow \textproc{rollout}(\pi_\theta, s_0)$ for $T=15$ steps
  \tcp*{sample $a_t\sim\mathcal{N}(a_\mu(s_t),\sigma^2)$}
$\hat A,\hat R \leftarrow \textproc{GAE}(\tau;\gamma{=}0.99,\lambda{=}0.95)$\;
$g_{\mathrm{PPO}} \leftarrow \nabla_\theta\big[-\mathcal{L}^{\mathrm{PPO}}(\theta;\tau,\hat A)+0.5\,\mathcal{L}^{V}(\phi;\tau,\hat R)-0.001\,\mathcal{H}(\pi_\theta)\big]$\;
\tcp{analytic branch: deterministic, differentiable, short horizon}
$\Lan(\theta) \leftarrow -\sum_{t=1}^{H} r\big(f_{\mathrm{sim}}^{\,t}(s_0,a_\mu(\cdot\,;\theta))\big)$\;
$\gan \leftarrow \nabla_\theta \Lan(\theta)$ \tcp*{\texttt{jax.grad} through the solver}
$g \leftarrow \beta\,\gan + (1-\beta)\,g_{\mathrm{PPO}}$\;
$\theta \leftarrow \textproc{Adam}(\theta, g;\ \eta{=}3\!\times\!10^{-4})$ \tcp*{single full-batch step}
\KwOut{$\theta,\phi$; \; logs $\big(\lVert\gan\rVert,\ \mathrm{RMS}\lvert a_t\rvert,\ \max\lvert a_t\rvert,\ \sum_t r_t\big)$}
PPO is this algorithm with the analytic branch (the $\Lan$, $\gan$, and mixing steps) removed, so that $g \leftarrow g_{\mathrm{PPO}}$.
\end{algorithm}

GPPO is the on-policy arm of the $2\times 2$
design and the direct object of the reproduction question this paper asks:
does the reported sample-efficiency gain hold, under the same fixed recipe,
at a chaotic operating point where the analytic gradient's premises ---
smooth, low-variance backpropagation through the solver --- are not
expected to hold? \cref{sec:results} answers this at $\Re=45$ (yes) and
$\Re=300$ (no seed-robust effect).

\subsection{GSAC: off-policy variant}
\label{sec:gsac}

GSAC applies the same analytic-gradient mix to the SAC actor rather than to
PPO's surrogate: after each episode's rollout and ten off-policy
critic-and-actor updates, one further actor step mixes the last actor
gradient, $g_{\mathrm{mf}}=g_\pi=\nabla_\varphi J_\pi(\varphi)$ from
\cref{eq:sacactor}, with a deterministic $H=6$-step analytic rollout from
the reset state (\cref{fig:gsac,alg:gsac}). The critics, target networks,
and replay mechanics are otherwise unchanged from plain SAC (\cref{sec:sac}).

This follows SAPO's core motivation of combining an off-policy,
maximum-entropy actor with an analytic gradient~\cite{sapo2025}, but
targets HydroGym's Kolmogorov flow rather than SAPO's rigid/deformable-body
platform, and uses the same fixed 50/50 mix as GPPO (\cref{sec:analytic})
rather than SAPO's entropy-based smoothing. To our knowledge this is the
first application of an off-policy, maximum-entropy analytic-gradient
method to a fluid-dynamics control task.

\begin{figure*}[t]
\centering
\begin{tikzpicture}[node distance=6mm and 9mm]
  \node[block] (roll) {rollout $T{=}15$\\ $\pi_\varphi$};
  \node[databox,right=of roll] (buf) {replay buffer $\mathcal{D}$\\ cap.\ $500\times 15$};
  \node[block,right=of buf] (batch) {sample\\ batch 64};
  \node[block,above=8mm of batch] (crit) {twin critics\\ $\psi_1,\psi_2$ \eqref{eq:saccritic}\\ $+$ targets $\bar\psi$};
  \node[block,right=of batch] (act) {actor grad $g_\pi$\\ \eqref{eq:sacactor}};
  \node[draw,dashed,rounded corners,fit=(batch)(crit)(act),inner sep=6pt] (inner) {};
  \node[font=\scriptsize,anchor=north] at ($(inner.south)+(0,-1pt)$) {$\times\,10$ updates / episode};
  \node[anblock,below=14mm of act] (an) {analytic branch\\ on actor, $H{=}6$\\ $\gan$ \eqref{eq:analytic}};
  \node[block,right=14mm of act,fill=yellow!12] (mix) {mix \eqref{eq:mix}\\ $\varphi\!\leftarrow$ Adam};
  \coordinate (top) at ($(inner.north)+(0,7mm)$);
  \draw[flow] (roll)--(buf); \draw[flow] (buf)--(batch);
  \draw[flow] (batch)--(crit); \draw[flow] (crit.east) -| (act.north);
  \draw[flow] (act.east) -- (mix.west);
  \draw[grad] (an.east) -| (mix.south);
  \draw[flow] (buf.south) |- (an.west);
  \draw[flow] (mix.north) -- (mix.north |- top) -- (roll.north |- top) -- (roll.north);
\end{tikzpicture}
\caption{GSAC training loop (one episode). Ten off-policy critic$+$actor
updates on replay batches are followed by a single analytic actor step
(6-step deterministic unroll from the reset state). Removing the red branch
and the mixer gives SAC.}
\label{fig:gsac}
\end{figure*}

\begin{algorithm}[t]
\caption{GSAC --- one training episode ($\beta=0.5$, $H=6$)}
\label{alg:gsac}
\DontPrintSemicolon
\KwIn{actor $\varphi$, critics $\psi_1,\psi_2$, targets $\bar\psi_1,\bar\psi_2$, replay buffer $\mathcal{D}$, reset state $s_0$}
$\tau \leftarrow \textproc{rollout}(\pi_\varphi, s_0)$ for $T=15$ steps;\quad $\mathcal{D} \leftarrow \mathcal{D}\cup\tau$\;
\For{$k = 1 \ldots 10$}{
  $B \leftarrow \textproc{sample}(\mathcal{D}, 64)$\;
  \lFor{$i\in\{1,2\}$}{$\psi_i \leftarrow \textproc{Adam}\big(\psi_i, \nabla_{\psi_i} J_Q(\psi_i; B)\big)$}
  $g_\pi \leftarrow \nabla_\varphi\,\mathbb{E}_{s\in B}\big[\alpha\log\pi_\varphi(a\mid s)-\min_i Q_{\psi_i}(s,a)\big],\ a\sim\pi_\varphi(\cdot\mid s),\ \alpha=0.1$\;
  $\varphi \leftarrow \textproc{Adam}(\varphi, g_\pi)$\;
  $\bar\psi_i \leftarrow (1-\tau_s)\bar\psi_i + \tau_s\psi_i,\quad \tau_s=0.005$\;
}
\tcp{analytic branch on the actor, after the SAC updates}
$\gan \leftarrow \nabla_\varphi\big[-\sum_{t=1}^{H} r\big(f_{\mathrm{sim}}^{\,t}(s_0,a_\mu(\cdot\,;\varphi))\big)\big]$\;
$\varphi \leftarrow \textproc{Adam}\big(\varphi,\ \beta\,\gan + (1-\beta)\,g_\pi\big)$ \tcp*{$g_\pi$ from the last inner step}
\KwOut{$\varphi,\psi$; \; logs $\big(\lVert\gan\rVert,\ \mathrm{RMS}\lvert a_t\rvert,\ \max\lvert a_t\rvert,\ \sum_t r_t\big)$}
SAC is this algorithm with the analytic branch (the final $\gan$ and mixing steps) removed.
\end{algorithm}

GSAC completes the $2\times 2$ design
(\cref{fig:design}) and lets this study ask its second question directly:
does an off-policy replay buffer tolerate the analytic gradient's
chaotic-regime instability better than an on-policy rollout, or does adding
a replay buffer simply add its own failure mode? Because GPPO and GSAC
share the identical analytic branch (\cref{eq:analytic}) and mixing rule
(\cref{eq:mix}), any difference in chaotic-regime robustness between them is
attributable to the on-/off-policy factor rather than to a different
analytic-gradient implementation.

GAE, PPO clipping, and the 50/50
analytic/model-free mix are hand-rolled in plain JAX/optax and kept
identical in form across the PPO and SAC training scripts. Both scripts
additionally log per-episode action amplitude (RMS and absolute max) for the
trigger-episode diagnostic of \cref{sec:mechanism}.
